% !TeX root = ../../Book1.tex
\section{Egoroff 定理}
\label{b1:sec:0405}

逐点收敛允许收敛开始的时刻随空间点改变，一致收敛则要求同一个时刻对所有点有效。两者通常相差很远。有限测度使这种差距可以集中到任意小的例外集上：舍去该集合以后，逐点收敛能够提升为一致收敛。

\subsection{逐点与一致收敛}
\label{b1:subsec:040501}

\cref{b1:def:pointwise-convergence}已经定义逐点收敛。对实值或复值函数列，它表示：对每个 \(x\in E\) 和每个 \(\eta>0\)，都存在可以依赖于 \(x\) 与 \(\eta\) 的整数 \(N=N(x,\eta)\)，使
\[
n\geq N
\quad\Longrightarrow\quad
|f_n(x)-f(x)|<\eta,
\]
也就是先固定空间点，再考察相应的数列。

\begin{definition}
\label{b1:def:uniform-convergence}
设 \(E\) 是集合，\(f_n,f:E\rightarrow\mathbb R\) 或
\(f_n,f:E\rightarrow\mathbb C\)。
如果对每个 \(\eta>0\)，存在只依赖于 \(\eta\) 的整数
\(N=N(\eta)\)，使
\[
n\geq N,\ x\in E
\quad\Longrightarrow\quad
|f_n(x)-f(x)|<\eta,
\]
就称 \(f_n\) 在 \(E\) 上\term[yizhi-shoulian]{一致收敛}[uniform convergence]于 \(f\)。
\end{definition}

两个定义的差别在量词次序。逐点收敛先固定 \(x\)，再选择 \(N\)；一致收敛先选择 \(N\)，这个 \(N\) 随后要适用于每个 \(x\)。当 \(E\ne\varnothing\) 时，后一条件等价于
\[
\sup_{x\in E}|f_n(x)-f(x)|\longrightarrow0,
\]
其中上确界允许暂时取 \(+\infty\)。一致收敛必然推出逐点收敛，反命题一般不成立。

回到第\ref{b1:subsec:040203}小节的 \(f_n(x)=x^n\)：它在 \([0,1]\) 上逐点收敛于 \(f=\mathbf 1_{\{1\}}\)，但
\[
\sup_{x\in[0,1]}|f_n(x)-f(x)|=1
\]
对每个 \(n\) 都成立，因为 \(x^n\) 在 \(x<1\) 时可以任意接近
\(1\)。所以收敛不一致。这个例子也说明，连续函数的逐点极限可以不连续。

\begin{proposition}
\label{b1:prop:uniform-limit-continuous}
若 \(E\subseteq\mathbb R^d\)，每个 \(f_n:E\rightarrow\mathbb R\) 都相对于 \(E\) 连续，且 \(f_n\to f\) 一致收敛，则 \(f\) 相对于 \(E\) 连续。复值函数同样成立。
\end{proposition}
\begin{proof}
固定 \(x\in E\) 与 \(\eta>0\)。由一致收敛，可取 \(N\)，使
\[
|f_N(y)-f(y)|<\frac{\eta}{3}
\quad(y\in E).
\]
由 \(f_N\) 在 \(x\) 处连续，存在 \(\delta>0\)，使得
\(y\in E\)、\(|x-y|<\delta\) 时
\[
|f_N(y)-f_N(x)|<\frac{\eta}{3}.
\]
于是
\begin{align*}
|f(y)-f(x)|
&\leq |f(y)-f_N(y)|+|f_N(y)-f_N(x)|+|f_N(x)-f(x)|\\
&<\eta.
\end{align*}
这正是 \(f\) 在 \(x\) 处相对于 \(E\) 的连续性。
\end{proof}

一致收敛也可以不预先知道极限而加以判定。

\begin{proposition}
\label{b1:prop:uniform-cauchy-criterion}
函数列 \(f_n:E\rightarrow\mathbb R\) 或 \(\mathbb C\) 一致收敛，当且仅当它满足\term[yizhi-cauchy-tiaojian]{一致 Cauchy 条件}[uniform Cauchy condition]：对每个 \(\eta>0\)，存在 \(N\)，使得
\[
m,n\geq N,\ x\in E
\quad\Longrightarrow\quad
|f_m(x)-f_n(x)|<\eta.
\]
\end{proposition}
\begin{proof}
若 \(f_n\to f\) 一致，取 \(N\) 使
\(|f_n(x)-f(x)|<\eta/2\) 对所有 \(n\geq N\)、\(x\in E\) 成立。三角不等式立即给出一致 Cauchy 条件。

反过来，设该条件成立。固定 \(x\in E\) 后，数列
\((f_n(x))\) 是 \(\mathbb R\) 或 \(\mathbb C\) 中的 Cauchy 列，由完备性收敛到某个值，记为 \(f(x)\)。给定 \(\eta>0\)，在一致 Cauchy 条件中把误差取为 \(\eta/2\)，得到 \(N\)。若 \(m\geq N\)，则对所有 \(n\geq N\) 和 \(x\in E\) 都有
\[
|f_m(x)-f_n(x)|<\frac{\eta}{2}.
\]
固定 \(m,x\) 并令 \(n\to\infty\)，得到
\[
|f_m(x)-f(x)|\leq\frac{\eta}{2}<\eta.
\]
同一个 \(N\) 对所有 \(x\) 有效，故 \(f_n\to f\) 一致收敛。
\end{proof}

\subsection{几乎一致收敛}
\label{b1:subsec:040502}

\begin{definition}
\label{b1:def:almost-uniform-convergence}
设 \((X,\mathcal A,\mu)\) 是测度空间，\(E\in\mathcal A\)，并设
\(f_n,f:E\rightarrow\mathbb R\) 或 \(\mathbb C\) 可测。如果对每个
\(\varepsilon>0\)，都存在可测集 \(A\subseteq E\)，使
\[
\mu(A)<\varepsilon
\]
并且 \(f_n\to f\) 在 \(E\setminus A\) 上一致收敛，就称
\(f_n\) 在 \(E\) 上\term[jihu-yizhi-shoulian]{几乎一致收敛}[almost uniform convergence]于 \(f\)。
\end{definition}

例外集 \(A\) 可以随 \(\varepsilon\) 改变。定义要求的是在整个
\(E\setminus A\) 上统一选择收敛指标，而不是对其中每个点分别选择。用上确界写，这个要求是
\[
\sup_{x\in E\setminus A}|f_n(x)-f(x)|
\longrightarrow0.
\]
若原收敛已经一致，可取 \(A=\varnothing\)，所以一致收敛蕴含几乎一致收敛。

把所有量词写出，几乎一致收敛的次序是
\[
\text{对每个 }\varepsilon>0\text{，存在 }A,\ \mu(A)<\varepsilon;
\quad
\text{对每个 }\eta>0\text{，存在 }N;
\quad
\text{对所有 }n\geq N,\ x\in E\setminus A.
\]
相比之下，几乎处处收敛只在开头删去一个零测集，随后逐点收敛所需的
\(N\) 仍可依赖于 \(x\)。Egoroff 定理的内容不是把例外集从正测度改成零测度，而是在有限测度条件下交换“选取 \(N\)”与“选取 \(x\)”的先后。

对于前面的 \(x^n\)，给定 \(\varepsilon>0\)，取 \(0<\delta<\min\{1,\varepsilon\}\) 并删去 \(A=(1-\delta,1]\)，便有 \(m(A)=\delta<\varepsilon\)，且
\[
\sup_{x\in[0,1-\delta]}|f_n(x)-f(x)|
\leq(1-\delta)^n\longrightarrow0.
\]
所以它几乎一致收敛。

这里不能把正测度但任意小的 \(A\) 换成一个固定零测集。事实上，若
\(N\subseteq[0,1]\) 为零测集，那么对每个 \(a<1\)，区间
\((a,1)\setminus N\) 都非空。即使把端点 \(1\) 放入 \(N\)，仍有
\[
\sup_{x\in[0,1]\setminus N}|x^n-f(x)|=1.
\]
所以 \(x^n\to f\) 不会在任何删去零测集后的定义域上一致收敛。几乎一致收敛允许例外集依精度改变，并允许它有正但任意小的测度。

另一个例子把例外集的作用显示得更直接。在 \(E=[0,1]\) 上令
\[
g_n=\mathbf 1_{[0,1/n]},
\quad
g=\mathbf 1_{\{0\}}.
\]
函数列逐点收敛，但不一致收敛，因为对每个 \(n\) 都能在
\((0,1/n]\) 中找到误差为 \(1\) 的点。给定 \(\varepsilon>0\)，取
\(0<\delta<\min\{1,\varepsilon\}\) 并删去 \(A=(0,\delta)\)。只要
\(n>\delta^{-1}\)，就在
\([0,1]\setminus A=\{0\}\cup[\delta,1]\) 上有 \(g_n=g\)。这里收敛甚至从某一项起完全相等。

\begin{proposition}
\label{b1:prop:almost-uniform-implies-ae}
在任意测度空间中，几乎一致收敛蕴含几乎处处收敛。
\end{proposition}
\begin{proof}
对每个 \(r\geq1\)，由几乎一致收敛取可测集
\(A_r\subseteq E\)，使 \(\mu(A_r)<2^{-r}\)，且
\(f_n\to f\) 在 \(E\setminus A_r\) 上一致收敛。令
\[
N=\bigcap_{r=1}^{\infty}A_r.
\]
对每个 \(r\)，都有 \(\mu(N)\leq\mu(A_r)<2^{-r}\)，故
\(\mu(N)=0\)。若 \(x\in E\setminus N\)，则至少存在一个 \(r\)，使
\(x\notin A_r\)；在 \(E\setminus A_r\) 上的一致收敛特别给出
\(f_n(x)\to f(x)\)。所以收敛在 \(E\setminus N\) 上逐点成立。
\end{proof}

这个方向不需要 \(\mu(E)<\infty\)。反向推理要把每个误差尺度下收敛过慢的点装进小集合，有限测度由此进入证明。

\subsection{Egoroff 定理的证明}
\label{b1:subsec:040503}

\term*[egoroff-dingli]{Egoroff 定理}[Egoroff theorem]把逐点收敛中的无穷多个点依赖指标集中到一个小例外集。先固定误差尺度，再控制从某一项起仍然违反该尺度的点。

\begin{theorem}
\label{b1:thm:egoroff}
设 \((X,\mathcal A,\mu)\) 是测度空间，\(E\in\mathcal A\) 且
\(\mu(E)<\infty\)。设可测函数
\(f_n,f:E\rightarrow\mathbb R\) 或 \(\mathbb C\) 在 \(E\) 上满足
\(f_n(x)\to f(x)\)。则 \(f_n\to f\) 在 \(E\) 上几乎一致收敛。
\end{theorem}
\begin{proof}
固定 \(\varepsilon>0\)。对整数 \(r,N\geq1\)，令
\[
B_{r,N}
=\bigcup_{n=N}^{\infty}
\{\,x\in E\mid |f_n(x)-f(x)|\geq r^{-1}\,\}.
\]
函数均可测，所以每个 \(B_{r,N}\) 可测。固定 \(r\) 后，集合列
\((B_{r,N})_N\) 随 \(N\) 递减。逐点收敛说明
\[
\bigcap_{N=1}^{\infty}B_{r,N}=\varnothing.
\]
一个点若属于这个交集，就会对每个 \(N\) 找到某个 \(n\geq N\)，使误差至少为 \(r^{-1}\)，这与该点处的收敛矛盾。

可测性在此处不能省略：它保证每一个误差集以及可数并
\(B_{r,N}\) 都可测，因而能够比较其测度并使用从上连续性。

又因 \(B_{r,1}\subseteq E\) 且 \(\mu(E)<\infty\)，测度的从上连续性给出
\[
\mu(B_{r,N})\longrightarrow0
\quad(N\to\infty).
\]
也可以改看好点集合
\[
G_{r,N}=E\setminus B_{r,N}
=\bigcap_{n=N}^{\infty}
\{\,x\in E\mid |f_n(x)-f(x)|<r^{-1}\,\}.
\]
对固定的 \(r\)，这些集合随 \(N\) 递增并覆盖 \(E\)。有限测度允许把
\(\mu(G_{r,N})\uparrow\mu(E)\) 改写成
\(\mu(E\setminus G_{r,N})\downarrow0\)。后一种写法就是坏点集合的估计。
因此可以为每个 \(r\) 选取整数 \(N_r\)，使
\[
\mu(B_{r,N_r})<\varepsilon2^{-r-1}.
\]
把所有误差尺度下的坏点放在一起，令
\[
A=\bigcup_{r=1}^{\infty}B_{r,N_r}.
\]
于是 \(A\subseteq E\) 可测，并且
\[
\mu(A)
\leq\sum_{r=1}^{\infty}\mu(B_{r,N_r})
\leq\frac{\varepsilon}{2}<\varepsilon.
\]

若 \(x\in E\setminus A\)，则对每个 \(r\) 都有
\(x\notin B_{r,N_r}\)。按照 \(B_{r,N_r}\) 的定义，只要
\(n\geq N_r\)，就有
\[
|f_n(x)-f(x)|<r^{-1}.
\]
这个 \(N_r\) 与 \(x\) 无关。给定 \(\eta>0\)，选择 \(r\) 使
\(r^{-1}<\eta\)，便知对所有 \(n\geq N_r\) 和所有
\(x\in E\setminus A\)，都有
\[
|f_n(x)-f(x)|<\eta.
\]
所以 \(f_n\to f\) 在 \(E\setminus A\) 上一致收敛。
\end{proof}

若只假设每个 \(f_n\) 可测，极限函数 \(f\) 的可测性也会由逐点极限保持可测性得到；定理把它列入假设，是为了让所有误差集合的可测性在陈述中直接可见。整个证明没有使用 \(X\) 上的拓扑，也不要求函数连续。与 Lusin 定理不同，Egoroff 定理是纯测度论结论。

证明可以直接按量词阅读。逐点收敛为每个 \(x\) 和每个精度
\(r^{-1}\) 提供一个依赖于 \(x\) 的起始位置；集合 \(B_{r,N}\) 记录起始位置晚于 \(N\) 的点。有限测度使这类点的测度随 \(N\) 趋于零。选出的 \(N_r\) 只依赖精度 \(r^{-1}\)，删除所有 \(B_{r,N_r}\) 后，便恢复了一致收敛所需的量词次序。

定理只要求 \(\mu(E)<\infty\)，不要求 \(E\) 有界、闭或紧。例如欧氏空间中可以取一个向无穷处延伸但总测度有限的可测集；只要它的测度有限，同一坏点集合论证仍然成立。紧集只在后面的 Lebesgue 正则性加强版中出现。

结合\cref{b1:prop:almost-uniform-implies-ae}可知，在
\(\mu(E)<\infty\) 时，几乎处处收敛与几乎一致收敛等价。前者允许先舍去一个零测集，后者则要求每个给定的正误差预算都留下一个一致收敛的大集合。下面把零测例外并入 Egoroff 定理。

\begin{corollary}[几乎处处版本]
\label{b1:cor:egoroff-ae}
设 \(\mu(E)<\infty\)，可测函数
\(f_n,f:E\rightarrow\mathbb R\) 或 \(\mathbb C\) 满足
\(f_n\to f\) 几乎处处。则 \(f_n\to f\) 几乎一致。
\end{corollary}
\begin{proof}
取可测零测集 \(N\subseteq E\)，使收敛在 \(E\setminus N\) 上逐点成立。给定 \(\varepsilon>0\)，把%
\cref{b1:thm:egoroff}应用于有限测度集 \(E\setminus N\)，得到可测集 \(A_0\subseteq E\setminus N\)，满足
\(\mu(A_0)<\varepsilon\)，且收敛在
\((E\setminus N)\setminus A_0\) 上一致。令 \(A=N\cup A_0\)。由于
\(\mu(N)=0\)，有 \(\mu(A)<\varepsilon\)，而
\[
E\setminus A=(E\setminus N)\setminus A_0,
\]
结论成立。
\end{proof}

同一结论适用于几乎处处有限的扩展实值函数列。若每个
\(f_n:E\rightarrow\overline{\mathbb R}\) 及极限 \(f\) 都几乎处处有限，并且
\(f_n\to f\) 几乎处处，那么先删去收敛失败的零测集，再删去
\[
\{\,|f|=\infty\,\}
\cup\bigcup_{n=1}^{\infty}\{\,|f_n|=\infty\,\}.
\]
这是零测集的可数并。余下集合上所有函数均为实值，可以使用上述推论。“几乎处处有限”这一条件使通常的绝对值误差在保留下来的集合上有定义。

在欧氏空间中还可把一致收敛集合取得紧一些。

\begin{corollary}
\label{b1:cor:egoroff-compact}
设 \(E\in\mathcal L^d\)、\(m(E)<\infty\)，且可测函数
\(f_n,f:E\rightarrow\mathbb R\) 或 \(\mathbb C\) 几乎处处收敛。则对每个
\(\varepsilon>0\)，存在紧集 \(K\subseteq E\)，使
\(m(E\setminus K)<\varepsilon\)，并且 \(f_n\to f\) 在 \(K\) 上一致收敛。
\end{corollary}
\begin{proof}
先由\cref{b1:cor:egoroff-ae}取可测集 \(A\subseteq E\)，使
\(m(A)<\varepsilon/2\)，且收敛在 \(F=E\setminus A\) 上一致。由内正则性在 \(F\) 中取紧集 \(K\)，使
\(m(F\setminus K)<\varepsilon/2\)。由于
\[
E\setminus K=A\sqcup(F\setminus K),
\]
所以 \(m(E\setminus K)<\varepsilon\)，而一致收敛在子集
\(K\subseteq F\) 上仍然成立。
\end{proof}

若上述推论中的每个 \(f_n\) 都相对于 \(E\) 连续，则限制
\(f_n|_K\) 在 \(K\) 上连续。结合%
\cref{b1:prop:uniform-limit-continuous}，可知极限 \(f|_K\) 也连续。这个结论常用于已有连续逼近列的场合；它与上一节 Lusin 定理的直接证明相容，但不能反过来充当那一证明，因为构造连续逼近列本身还需要额外论证。

有限测度假设不能删除。令 \(E=\mathbb R\)，并取
\[
f_n=\mathbf 1_{[n,n+1)},
\quad
f=0.
\]
对每个固定的 \(x\in\mathbb R\)，只有有限多个区间 \([n,n+1)\) 能含
\(x\)，所以 \(f_n(x)\to0\)。但若可测集 \(A\subseteq\mathbb R\) 满足
\(m(A)<1/2\)，那么对每个 \(n\) 都有
\[
m\bigl([n,n+1)\setminus A\bigr)
\geq1-m(A)>\frac12.
\]
故 \([n,n+1)\setminus A\) 非空，并且
\[
\sup_{x\in\mathbb R\setminus A}|f_n(x)|=1
\quad(n\geq1).
\]
收敛在 \(\mathbb R\setminus A\) 上不一致，因此也不几乎一致。在定理的证明中，有限测度正是用来从
\(B_{r,N}\downarrow\varnothing\) 推出 \(\mu(B_{r,N})\to0\)；本例在
\(r=2\) 时有 \(B_{2,N}=[N,\infty)\)，其测度始终为无穷。

无穷测度空间上的局部结论仍然成立。若
\(E=\bigcup_{j=1}^{\infty}E_j\)，其中 \(E_j\uparrow E\) 且
\(\mu(E_j)<\infty\)，则可在每个 \(E_j\) 上分别使用 Egoroff 定理。不过，所得一致收敛指标通常依赖于 \(j\)，不能拼成整个 \(E\) 上的一致收敛；移动区间的例子正表明这种局部结论不能自动升级为全局结论。

至此，Lusin 定理与 Egoroff 定理的分工也清楚了：前者用欧氏正则性把可测函数改造成大集合上的连续函数，后者只用有限测度把逐点收敛加强为几乎一致收敛。Walter Rudin 的《Real and Complex Analysis》\cite{Rudin1987Real}给出了较紧凑的对照叙述；阅读时仍要把这两组假设分开核对。
